% !TeX spellcheck = hu_HU
% !TeX encoding = UTF-8
% !TeX program = xelatex
%----------------------------------------------------------------------------
\chapter{A blAST eszköz}
%----------------------------------------------------------------------------

Ebben a fejezetben mutatom be a saját fejlesztésű blAST eszközt.
Ennek az eszköznek a célja, hogy egy magas szintű, dinamikus nyelvben biztosítson lehetőséget projekt specifikus ellenőrzések megvalósítására.

\section{A C4 nyelv}

A használt nyelv a saját fejlesztésű C4 programozási nyelv \cite{bodorAltalanosCeluProgramnyelv2025}.
Ez egy immutábilis, dinamikusan típusos, lusta kiértékelésű funkcionális nyelv, amelyik az LLVM infrastruktúrát felhasználva natív kódra fordul.

A nyelv erőssége, hogy direkt céljának tekinti, hogy beágyazott szakterület-specifikus nyelveket (\gls{eDSL}) lehessen benne könnyedén magvalósítani.
Ennek okán, nagyon minimális szintaktikával rendelkezik, és a függvények kompozíciójából állítja elő még a gyakori nyelvi elemeket is.
Ez a rugalmasság biztosítja, hogy a blAST nevű eszköz nagyon szabadon szabhatja meg a számára előnyős nyelvi elemeket.

A nyelv részletes tárgyalása most nem fog megtörténni, csak a fontos nyelvi elemeket mutatom be, hogy az alábbiak érhetőek legyenek.
A megvalósítás belső részeiről a továbbiakban szó esik, ahol azon részletek fontossá válnak, itt magáról a nyelvről egy általános leírás szerepel.

\subsection*{Függvények}

A C4 függvények létrehozásához a {\tt\bf let} kulcsszót lehet használni, amelyik utána a függvény nevét, aritását, majd törzsét kell megadni.
Mivel dinamikusan típusos a nyelv, így típusokat nem kell feltüntetni, az artiásnak azonban egyeznie kell.
Deklarálni a törzs helyett egy pontosvessző megadásával lehet.

A változók is függvények, így ugyanezen szintaktika tartozik definiálásukhoz.
Annyi a különbség, hogy nulla paramétert vesznek át (0 az aritásuk), és egy nem blokk kifejezést tartalmaznak.

Függvényt meghívni, csak ha az előre ismert lehetséges.
Ilyenkor, a függvény neve után bármilyen szóközzel (whitespace) elválasztva kell megadni 

\begin{lstlisting}
let foo/2 { |a b| a + b }
io_println foo 1 2
\end{lstlisting}

Változókban tárolt vagy függvényből visszakapott kifejezéseket függvényként meghívni a {\tt{\bf \&({\sl EXPR})/{\sl N}} ARG...} struktúrával lehet.
Az {\sl EXPR} kifejezés helyén bármilyen kifejezés állhat -- még primitív értékeket, mint az egészeket is meg lehet hívni.
A paramétereket ilyen esetben nem tudja a fordító kitalálni, így meg kell adni hány darab {\sl ARG} kifejezés áll argumentumként.
Többet lehet megadni, ezek pusztán nem kerülnek felhasználásra, kevesebb esetén a működés definiálatlan.
Elsődleges felhasználás blokk értékek (más nyelvekben: lambdák, függvény objektumok, delegate-ek) meghívása.

\begin{lstlisting}
let make_adder/1 { |x| 
	{ |y| x + y } # ez a blokk a visszatérési érték
}
io_println &(42)/1 "szöveg"      # => 42 (fölösleges paraméter figyelmen kívül hagyva)
io_println &(make_adder 20)/1 22 # => 42
\end{lstlisting}

\subsection*{Operátorok}

Operátorokat a függvényekhez hasonlóan lehet bevezetni.
Teljesen saját operatárokat is lehet készíteni, azonban a komment karakterrel (kettőskereszt (\#)) nem kezdődhetnek.
Bináris operátor esetén a precedenciát és asszociativitást is meg kell adni, ami alapján fognak elemzés során zárójeleződni.

\begin{lstlisting}
let (~)/1 { |x| -x }
let (~)/2 left 30 { |x y| x - y }
io_println 1 ~ 2 ~ ~1 # = ((1 - 2) - (-1))
\end{lstlisting}

\subsection*{Csomagok}

A C4 nyelv csomagokba való szervezést statikus könyvtárak (archívumok) segítségével teszi lehetővé.
Minden függvény (és operátor) lehet publikus ($+$), belső használatú (internal, \textasciitilde) és privát ($-$).
A publikus függvények láthatóak lesznek a könyvtár interfészén is, a belső használatúak csak másik lokális fájlokban elérhetőek.
Privátok mindig csak a saját fájlban láthatóak.

Alapból minden belső használatú, minden más láthatóságot a függvény neve előtt kell megadni a {\tt\bf let} kulcsszó után.

Használni másik fájlokat kétféle módon lehet: a {\tt {\bf use} "file.c4"} szintaktika egy lokális fájl függvényeit teszi elérhetővé, 
a {\tt {\bf use} <lib.c4a>} pedig egy könyvtárét.

\subsection*{Natív interoperabilitás}

C és C4 között támogatott az idegen függvények (foreign function) hívása.
A fordító a konkrét lustán kiértékelt környezetből való natív függvény hívásához szükséges munka nagy részét elvégzi.

A C4 oldalon a függvényt csak deklarálni szabad, és a {\tt\bf native} kontextus függő kulcsszót kell használni.
Ilyenkor egy ``c4sym\_''-el prefixelt nevű C függvényt fog keresni a linker a program összeállításakor.
Ennek annyi paramétere kell, hogy legyen, amennyi az aritásként lett megadva a C4 függvény deklarálásakor, különben a viselkedés nem definiált.
Indirekt meghíváskor ennek is lehet azonban több paramétert átadni.

A C-ben definiált függvény minden paramétere {\tt c4\_datum}, amelyik kezeléséhez függvények a ``c4rt3/c4rt.h'' fájlban találhatóak.
A konkrét definiáláskor használható az ugyanezen fejlécben található {\tt {\bf c4\_let\_native}()} segédmakró, aminek a konkrét C4 oldali függvénynevet elég megadni.

\begin{lstlisting}
let native foo/2;
# c4_let_native(foo)(c4_datum a, c4_datum b) { ... }
\end{lstlisting}

A C oldalról a {\tt c4\_call} és {\tt c4\_apply} függvényekkel lehetséges egy objektum kiértékelését kérni (a nyelv lusta volta miatt lehetnek kiértékeletlen értékek) illetve egy kifejezést átadott paraméterekkel meghívni.
Utóbbi analóg a {\tt \&()/N} szintaktikával.

Érdemes megjegyezni, hogy a C oldalon nem garantálható az immutábilitás: a paraméterekben kapott értéket C kód számára módosíthatóak.
Ez azonban nem jól viselkedő programokat eredményezhet, így hivatalosan nem definiált működést jelent.
Amit azonban lehetséges szabadon módosítani, azok egyéb C kódban elérhető globális változók.
Ezekre nem tesz megkötést a rendszer, és emiatt azt sem feltételezi, hogy natív függvények függvényként viselkednek: függhetnek tehát a C globálisok állapotától is.

\input content/blast-prime

\input content/blast-arch

